\documentclass[11pt,letterpaper]{article}
\usepackage{cornell-notes}
\usepackage{needspace}

\newcommand{\CornellDocumentTitle}{Chapter 102: ai abuse investigating the threat landscape}
\title{\CornellDocumentTitle}
\author{}
\date{}
\setDocTitle{\CornellDocumentTitle}
\setDocSubtitle{Cybersecurity Cornell Notes}
\setCornellCollection{Cybersecurity Reference Notes}
\setCornellUnitType{Chapter}
\setCornellUnitNumber{102}
\setCornellUnitTitle{ai abuse investigating the threat landscape}
\setCornellCompanionLabel{Cybersecurity study companion}

\begin{document}
\maketitle
\begin{CornellOverviewBox}{Chapter Orientation}
\textbf{Chapter orientation.} This chapter examines ai abuse: investigating the threat landscape within the broader domain of advanced security. Its central problem is how to reason about accuracy, classifier, attacks, and training while keeping security objectives, operating assumptions, evidence, and recovery connected. A practitioner should approach the material through capability, data and model provenance, misuse paths, uncertainty, human oversight, and defense under changing behavior. Useful prerequisites include basic risk, threat, control, and systems concepts.
\end{CornellOverviewBox}
\section*{Learning Objectives}
\begin{itemize}
\item Explain the security purpose and scope of AI Abuse: Investigating the Threat Landscape.
\item Identify the assets, actors, assumptions, and trust boundaries associated with accuracy, classifier, and attacks.
\item Distinguish Introduction from Background and explain where their responsibilities intersect.
\item Analyze how failures in accuracy, classifier, and attacks can affect confidentiality, integrity, availability, privacy, or safety.
\item Evaluate relevant controls through capability, data and model provenance, misuse paths, uncertainty, human oversight, and defense under changing behavior.
\item Audit an implementation using model and data documentation, abuse cases, evaluations, access records, monitoring, incident evidence, and governance decisions.
\item Prioritize remediation according to exploitability, consequence, control strength, and recovery readiness.
\item Verify that corrective actions change observable risk rather than documentation alone.
\end{itemize}
\section*{Cornell Cue-and-Notes}
\Needspace{8\baselineskip}
\subsection*{Introduction}
\begin{longtable}{>{\RaggedRight\arraybackslash}p{0.29\textwidth}|>{\RaggedRight\arraybackslash}p{0.65\textwidth}}
\CornellHeader
\CornellNoteRow{What security problem does Introduction address?}{\textbf{Core idea.} Treat Introduction as a structured part of AI Abuse: Investigating the Threat Landscape, with particular attention to learning, attacks, adversarial, and malicious. \textbf{Analytical lens.} This topic establishes a component of the chapter's argument; connect its definitions and mechanisms to a testable security objective. For learning, attacks, and adversarial, test misuse and evasion while keeping model or automation authority bounded and reversible. \textbf{Security reasoning.} Identify what must remain protected, who or what can alter the expected state, which assumptions cross a trust boundary, and how failure changes confidentiality, integrity, availability, privacy, or safety. \textbf{Application.} The practitioner should separate demonstrated behavior from forecast, test misuse cases, and retain accountable human control. \textbf{Evidence.} Seek model and data documentation, abuse cases, evaluations, access records, monitoring, incident evidence, and governance decisions. Related subtopics include Adversarial Attacks Against ML.}
\end{longtable}
\begin{CornellSummaryBox}{Topic Summary}
Introduction is best remembered as the connection between learning, attacks, adversarial, and malicious and the chapter's larger control objective. A defensible review states the assumption, expected behavior, evidence, failure consequence, and required response.
\end{CornellSummaryBox}
\Needspace{8\baselineskip}
\subsection*{Background}
\begin{longtable}{>{\RaggedRight\arraybackslash}p{0.29\textwidth}|>{\RaggedRight\arraybackslash}p{0.65\textwidth}}
\CornellHeader
\CornellNoteRow{Which assumptions matter in Background?}{\textbf{Core idea.} Treat Background as a structured part of AI Abuse: Investigating the Threat Landscape, with particular attention to attacks, label, training, and learning. \textbf{Analytical lens.} This topic establishes a component of the chapter's argument; connect its definitions and mechanisms to a testable security objective. For attacks, label, and training, test misuse and evasion while keeping model or automation authority bounded and reversible. \textbf{Security reasoning.} Identify what must remain protected, who or what can alter the expected state, which assumptions cross a trust boundary, and how failure changes confidentiality, integrity, availability, privacy, or safety. \textbf{Application.} The practitioner should separate demonstrated behavior from forecast, test misuse cases, and retain accountable human control. \textbf{Evidence.} Seek model and data documentation, abuse cases, evaluations, access records, monitoring, incident evidence, and governance decisions. Related subtopics include Defending Against Adversarial Attacks to ML, and Taxonomy of Attacks.}
\end{longtable}
\begin{CornellSummaryBox}{Topic Summary}
Background is best remembered as the connection between attacks, label, training, and learning and the chapter's larger control objective. A defensible review states the assumption, expected behavior, evidence, failure consequence, and required response.
\end{CornellSummaryBox}
\Needspace{8\baselineskip}
\subsection*{Attack Models}
\begin{longtable}{>{\RaggedRight\arraybackslash}p{0.29\textwidth}|>{\RaggedRight\arraybackslash}p{0.65\textwidth}}
\CornellHeader
\CornellNoteRow{How should a reviewer evaluate Attack Models?}{\textbf{Core idea.} Treat Attack Models as a structured part of AI Abuse: Investigating the Threat Landscape, with particular attention to learning, attack, machine, and samples. \textbf{Analytical lens.} This topic is threat-centered: separate actor capability, reachable attack surface, prerequisites, execution path, and consequence. For learning, attack, and machine, test misuse and evasion while keeping model or automation authority bounded and reversible. \textbf{Security reasoning.} Identify what must remain protected, who or what can alter the expected state, which assumptions cross a trust boundary, and how failure changes confidentiality, integrity, availability, privacy, or safety. \textbf{Application.} The practitioner should separate demonstrated behavior from forecast, test misuse cases, and retain accountable human control. \textbf{Evidence.} Seek model and data documentation, abuse cases, evaluations, access records, monitoring, incident evidence, and governance decisions. Related subtopics include Preparation, and Manifestation.}
\end{longtable}
\begin{CornellSummaryBox}{Topic Summary}
Attack Models is best remembered as the connection between learning, attack, machine, and samples and the chapter's larger control objective. A defensible review states the assumption, expected behavior, evidence, failure consequence, and required response.
\end{CornellSummaryBox}
\Needspace{8\baselineskip}
\subsection*{Case Studies}
\begin{longtable}{>{\RaggedRight\arraybackslash}p{0.29\textwidth}|>{\RaggedRight\arraybackslash}p{0.65\textwidth}}
\CornellHeader
\CornellNoteRow{Why can Case Studies fail in practice?}{\textbf{Core idea.} Treat Case Studies as a structured part of AI Abuse: Investigating the Threat Landscape, with particular attention to accuracy, classifier, dataset, and label. \textbf{Analytical lens.} This topic is evidence by example: extract the reusable mechanism while keeping environmental assumptions and limits visible. For accuracy, classifier, and dataset, test misuse and evasion while keeping model or automation authority bounded and reversible. \textbf{Security reasoning.} Identify what must remain protected, who or what can alter the expected state, which assumptions cross a trust boundary, and how failure changes confidentiality, integrity, availability, privacy, or safety. \textbf{Application.} The practitioner should separate demonstrated behavior from forecast, test misuse cases, and retain accountable human control. \textbf{Evidence.} Seek model and data documentation, abuse cases, evaluations, access records, monitoring, incident evidence, and governance decisions. Related subtopics include Adversarial Attacks Against ML-Based, Detection of Fake Twitter Accounts, Attack Evaluation, and Methodology.}
\end{longtable}
\begin{CornellSummaryBox}{Topic Summary}
Case Studies is best remembered as the connection between accuracy, classifier, dataset, and label and the chapter's larger control objective. A defensible review states the assumption, expected behavior, evidence, failure consequence, and required response.
\end{CornellSummaryBox}
\begin{CornellWarningBox}{Historical Context}
Some products, protocols, examples, threat observations, or legal references in this chapter are historically bounded. Retain the underlying threat model and control objective, but verify present applicability before treating a named implementation as current guidance.
\end{CornellWarningBox}
\begin{CornellOverviewBox}{Defensive Guidance}
Apply the chapter by making assets, authority, trust, expected behavior, and failure response explicit. In practice, separate demonstrated behavior from forecast, test misuse cases, and retain accountable human control. A control is not fully supported until its operation, monitoring, ownership, and recovery behavior are demonstrated with evidence.
\end{CornellOverviewBox}
\section*{Threat-and-Control Map}
\begingroup\scriptsize\setlength{\tabcolsep}{2pt}
\begin{longtable}{>{\RaggedRight\arraybackslash}p{0.135\textwidth}>{\RaggedRight\arraybackslash}p{0.145\textwidth}>{\RaggedRight\arraybackslash}p{0.155\textwidth}>{\RaggedRight\arraybackslash}p{0.145\textwidth}>{\RaggedRight\arraybackslash}p{0.16\textwidth}>{\RaggedRight\arraybackslash}p{0.17\textwidth}}
\toprule\rowcolor{Navy}\color{white}\textbf{Asset} & \color{white}\textbf{Threat / Failure} & \color{white}\textbf{Vulnerability} & \color{white}\textbf{Impact} & \color{white}\textbf{Detection / Evidence} & \color{white}\textbf{Control}\\\midrule
\endfirsthead
\toprule\rowcolor{Navy}\color{white}\textbf{Asset} & \color{white}\textbf{Threat / Failure} & \color{white}\textbf{Vulnerability} & \color{white}\textbf{Impact} & \color{white}\textbf{Detection / Evidence} & \color{white}\textbf{Control}\\\midrule
\endhead
Models, training data, automated decisions, and dependent users or systems & Abuse, evasion, poisoning, manipulation, leakage, or unsafe automation & Opaque behavior, weak provenance, excessive authority, or inadequate evaluation & Scaled error, privacy loss, compromised decisions, or accelerated attacks & Evaluation results, monitoring, provenance checks, abuse reporting, and independent review & Scoped authority, secure data and model lifecycles, adversarial testing, oversight, and rollback\\\midrule
Assumptions surrounding accuracy & Misuse or unexpected state transition & Trust in classifier is not validated & A local weakness becomes a broader compromise path & Review evidence for attacks and test negative paths & Make trust explicit, constrain authority, and fail safely\\\midrule
Recovery capability and decision evidence & Control failure remains unnoticed or cannot be reversed & Monitoring, ownership, or restoration has not been exercised & Longer exposure and slower, less reliable recovery & Alert-to-response exercises and restoration tests & Assigned ownership, actionable telemetry, preserved evidence, and rehearsed recovery\\\midrule
\bottomrule
\end{longtable}
\endgroup
\section*{Practical Security Checklist}
\begin{itemize}[label=$\square$]
\item Define the in-scope assets, users, services, data, and operational dependencies for AI Abuse: Investigating the Threat Landscape.
\item Document the expected behavior and security objective of accuracy.
\item Map identities, privileges, trust boundaries, and externally controlled inputs associated with accuracy and classifier.
\item Record the threat or failure being addressed, its prerequisites, likely path, and credible consequences.
\item Inspect model and data documentation, abuse cases, evaluations, access records, monitoring, incident evidence, and governance decisions.
\item Test both the intended path and negative paths, including invalid state, partial failure, excessive privilege, and unavailable dependencies.
\item Confirm preventive controls are complemented by actionable detection and a named response owner.
\item Exercise containment, restoration, and attacks; record actual recovery behavior and unresolved dependencies.
\item Validate exceptions, compensating controls, expiration dates, and accountable approval.
\item Preserve reproducible evidence and tie each finding to affected assets, risk, remediation, and verification criteria.
\end{itemize}
\begin{CornellSummaryBox}{Chapter Synthesis}
The chapter's principal lesson is that ai abuse: investigating the threat landscape must be evaluated as an operating security system rather than an isolated product or statement. The most important failure patterns involve accuracy, classifier, attacks, and training, especially when ownership, boundaries, telemetry, or recovery remain implicit. A reviewer should connect architecture and behavior to model and data documentation, abuse cases, evaluations, access records, monitoring, incident evidence, and governance decisions, prioritize findings by reachable consequence and control independence, and verify remediation through observable change. Within Advanced Security, this chapter contributes a reusable method: state the objective, expose assumptions, test failure, preserve evidence, and learn from operation.
\end{CornellSummaryBox}
\section*{Self-Test Questions}
\begin{enumerate}
\item What is the central security objective of AI Abuse: Investigating the Threat Landscape?
\item How does Introduction contribute to the chapter's broader security argument?
\item Distinguish Background from Attack Models.
\item Which trust assumptions should be made explicit when evaluating accuracy?
\item How could a weakness involving classifier affect confidentiality, integrity, availability, privacy, or safety?
\item What evidence would demonstrate that controls surrounding attacks operate as intended?
\item Why should prevention, detection, response, and recovery be evaluated together in this domain?
\item Scenario: A powerful automated capability is deployed for defensive benefit without testing how an adversary could redirect or evade it. What should the reviewer examine first, and why?
\item Scenario: A control associated with training passes a configuration check but fails during an operational exercise. How should the finding be interpreted and prioritized?
\item How should a practitioner distinguish enduring principles in this chapter from time-sensitive tools, examples, or implementation details?
\end{enumerate}
\Needspace{8\baselineskip}
\section*{Answer Key}
\begin{enumerate}
\item The objective is to protect the relevant assets by understanding capability, data and model provenance, misuse paths, uncertainty, human oversight, and defense under changing behavior and by connecting controls to measurable risk reduction.
\item Introduction provides one part of the causal chain: it should identify expected behavior, dependencies, failure modes, and the evidence needed to justify confidence.
\item Background and Attack Models address different portions of the problem. A strong comparison states their scope, assumptions, enforcement points, evidence, and how a failure in one affects the other.
\item Reviewers should identify who controls accuracy, what inputs or dependencies it trusts, where privilege changes, what happens during partial failure, and how those assumptions are tested.
\item A weakness in classifier can propagate across trust boundaries. The actual impact depends on reachable assets, granted authority, control independence, visibility, and recovery capability.
\item Useful evidence includes model and data documentation, abuse cases, evaluations, access records, monitoring, incident evidence, and governance decisions; configuration alone is insufficient when operation, monitoring, or recovery is part of the claim.
\item Prevention reduces probability, detection limits dwell time, response limits spread, and recovery restores objectives. Omitting one layer creates an unexamined failure mode.
\item Begin by establishing scope, ownership, expected behavior, and the violated assumption. Then separate demonstrated behavior from forecast, test misuse cases, and retain accountable human control, preserving evidence for prioritization and follow-up.
\item Treat the exercise as stronger evidence than a static check. Prioritize according to reachable impact, likelihood of real operating conditions, missing compensating controls, and recovery difficulty.
\item Retain the principle, threat model, and control objective; separately label technologies, legal references, product examples, and threat observations whose applicability depends on time or environment.
\end{enumerate}
\section*{Key-Term Recap}
\begin{longtable}{>{\RaggedRight\arraybackslash}p{0.27\textwidth}>{\RaggedRight\arraybackslash}p{0.66\textwidth}}
\toprule\rowcolor{Navy}\color{white}\textbf{Term} & \color{white}\textbf{Meaning}\\\midrule
\endfirsthead
\toprule\rowcolor{Navy}\color{white}\textbf{Term} & \color{white}\textbf{Meaning}\\\midrule
\endhead
\textbf{Threat} & A circumstance or actor capable of causing harm by exploiting a weakness or adverse condition. \\\midrule
\textbf{Intrusion Detection} & Monitoring and analysis intended to identify unauthorized or malicious activity. \\\midrule
\textbf{Malware} & Software or code intentionally designed to disrupt, damage, spy, or obtain unauthorized control. \\\midrule
\textbf{Threat Model} & A structured representation of assets, trust boundaries, adversaries, attack paths, and mitigations. \\\midrule
\textbf{Vulnerability} & A weakness that can be exploited or triggered to violate a security objective. \\\midrule
\textbf{Resilience} & The ability to anticipate, withstand, recover from, and adapt to adverse conditions. \\\midrule
\textbf{Integrity} & The property that information and systems remain accurate, complete, and protected from unauthorized modification. \\\midrule
\textbf{Risk} & The effect of uncertainty on objectives, commonly evaluated through likelihood, impact, exposure, and context. \\\midrule
\textbf{Artificial Intelligence} & Computational techniques that perform tasks such as classification, prediction, or generation and introduce data, model, and misuse risks. \\\midrule
\textbf{Attack Surface} & The collection of reachable interfaces, inputs, privileges, dependencies, and operational paths an adversary might exploit. \\\midrule
\bottomrule
\end{longtable}
\begin{CornellExamBox}{One-Sentence Takeaway}
Treat ai abuse: investigating the threat landscape as a verifiable chain from assets and assumptions through controls, evidence, response, and recovery.
\end{CornellExamBox}
\end{document}
